\documentclass[11pt]{article}
\usepackage[margin=25mm]{geometry}
\usepackage{fontspec}
\setmainfont{Latin Modern Roman}
\usepackage{amsmath,amssymb,mathtools}
\usepackage{unicode-math}
\setmathfont{Latin Modern Math}
\usepackage{xurl,hyperref}
\hypersetup{colorlinks=true,urlcolor=blue}
\setcounter{secnumdepth}{0}
\setlength{\parindent}{0pt}
\setlength{\parskip}{5pt}
\setlength{\emergencystretch}{3em}
\newcommand{\sref}[2]{\hyperlink{source:#1:#2}{[S#2]}}
\newcommand{\cataloguescope}[1]{\paragraph{Recorded target.}#1}
\newcommand{\R}{\mathbb R}
\newcommand{\C}{\mathbb C}
\newcommand{\N}{\mathbb N}
\newcommand{\Z}{\mathbb Z}
\newcommand{\Q}{\mathbb Q}
\newcommand{\F}{\mathbb F}
\newcommand{\eps}{\varepsilon}
\begin{document}
% UnsolvedMath ID 20002550; source catalogue code AIM-PHYSICS-0058
\setcounter{section}{20002549}
\section{Level-orbit potentials: quantifier audit, speed criterion, and transnormal rigidity}\label{20002550}
{\small UnsolvedMath ID 20002550 · Mathematical Physics}
\subsection{Definitions and mathematical statement}

\paragraph{Shared notation.}$\mathbb N=\{1,2,\ldots\}$; the standard number fields have their usual meanings. All additional conventions are stated locally.


\paragraph{Definitions and notation.}For a planar conservative system a smooth potential $U$ gives motion $\ddot x=-\nabla U(x)$ under the physical-potential convention. Tangency to a regular level at initial time means $\nabla U(x(0))\cdot\dot x(0)=0$. Remaining on that level requires $U(x(t))$ to be constant throughout the motion. The speed as well as the initial direction affects the curvature needed to stay on a level.

\paragraph{Catalogue wording retained for traceability.}
Does a planar potential $U(x)$ exist with the property that if $\dot x(0)$ is tangent to a level curve of $U$, then $U(x(t))$ remains constant?
\paragraph{Scope and source qualification.}The source does not exclude constant potentials, specify regularity or say whether every tangential speed or a selected speed is meant. These choices change the answer and must be identified; no nontrivial universal quantifier is invented. The inherited solved label is not a fresh proof endorsement. The catalogue title/status was paired with a machine-generated research attempt; that attempt is not peer-reviewed evidence of settlement. The statement covers the original source questions rather than only the title’s subcase.
\subsection{Short English statement}
Find planar potentials whose trajectories remain on an initial equipotential when launched tangentially.
\subsection{Research attempt}
\paragraph{Model, date and investigated scope.}
OpenAI GPT-6 Astra (configured model \texttt{gpt-6-astra}), ultra reasoning; generated 11 October 2026. Proved the every-tangential-speed obstruction and explicit selected-speed harmonic circles, distinguishing constant-potential existence. The intended nontrivial formulation and general selected-speed classification remain unspecified; source-solved status is preserved.

\paragraph{Sources and literature check.}
Read the supplied complete definition, catalogue wording, scope qualifications and source list for this canonical ID. The original-source citations are inherited as supplied references; no fresh primary-source or novelty search was performed for this attempt. No current open-status claim is inferred from this calculation.

\paragraph{The every-speed interpretation.}Let \(U\in C^2\) on a connected open planar domain and use \(\ddot x=-\nabla U\). Along any classical solution,
\[\frac{d^2}{dt^2}U(x(t))=\dot x^TD^2U\dot x-|\nabla U|^2.
\]
At a regular point \(\nabla U\ne0\), choose a unit tangent \(e\) and launch with velocity \(se\). Remaining on the level would require \(s^2e^TD^2Ue=|\nabla U|^2\). This cannot hold for two distinct positive values of \(s^2\): subtraction forces the Hessian term zero, then the gradient zero, a contradiction. Including zero velocity gives the contradiction immediately. Thus a property demanded for every tangential speed has no regular points; if imposed everywhere, \(\nabla U=0\) throughout and \(U\) is constant on the connected domain.
\paragraph{A selected-speed positive example.}For \(U(x)=\tfrac12\omega^2|x|^2\), \(\omega>0\), a point on a circle of radius \(r>0\) with tangential speed \(\omega r\) evolves as uniform circular harmonic motion, staying on that level. A different speed generally gives an ellipse on which \(|x|\) varies. Directly, with \(x(0)=(r,0)\), \(\dot x(0)=(0,s)\), the solution is \((r\cos\omega t,(s/\omega)\sin\omega t)\); its squared radius is constant exactly when \(s^2=\omega^2r^2\).
\paragraph{Source qualification.}Constant potentials answer the unrestricted existence wording trivially, whereas selected-speed and every-speed nonconstant questions differ. This reconstructs the quantifier distinction and keeps the inherited source-solved label without claiming a newly discovered settlement. A general selected-speed classification and the precise intended global domain remain outside the calculation. No change to the historical question is silently made.

\paragraph{Verification and adversarial audit.}
Checked chain-rule second derivative, two-speed contradiction, zero velocity, connected constant case and exact harmonic-oscillator ellipse test. This was an LLM adversarial audit, not independent human verification; no theorem prover was used. Prompts A then B applied to this ID: Checked chain-rule second derivative, two-speed contradiction, zero velocity, connected constant case and exact harmonic-oscillator ellipse test. Scope remains limited: Fix the intended speed quantifier and global regularity/domain; classify selected-speed level-preserving potentials if that is the desired nontrivial target.

Coordinator checked the second derivative along a potential level, the two-distinct-speeds obstruction and the isotropic-oscillator selected-speed example. This additional check was model-assisted, not independent human verification.

\paragraph{Tools, process and contributions.}
Python 3 and local shell used to read supplied files and serialize this manuscript; mathematics derived symbolically in the text, with no numerical experiment or formal proof assistant. Contribution prompt A was followed by prompt B. GPT-6 Astra with ultra reasoning drafted the research and performed its stated self-audit; any additional coordinator contribution is identified in the verification and audit. A GPT-6-based Codex coordinator assembled the manuscript and checked repository integration; its exact variant and reasoning configuration were not exposed. No human mathematical contribution, independent human proof check or formal proof certification is supplied. Complete derivations and the consulted source references are included here; private reasoning transcripts are not published.

\paragraph{Final status and first remaining gap.}
\textbf{UNRESOLVED}. Fix the intended speed quantifier and global regularity/domain; classify selected-speed level-preserving potentials if that is the desired nontrivial target.

\section{Completion Estimate}
\noindent\textbf{COMPLETION: 0\%}
The calculation is an elementary benchmark or conditional reduction; it establishes no claimed new portion of the complete source target.

\subsection{Sources}
\begin{itemize}
\item[\hypertarget{source:20002550:1}{S1}] ULAM.AI and the UnsolvedMath Contributors, UnsolvedMath, record 20002550 (AIM-PHYSICS-0058). Catalogue curation under CC BY 4.0; definitions and English exposition added locally.
\url{https://huggingface.co/datasets/ulamai/UnsolvedMath}.
\item[\hypertarget{source:20002550:2}{S2}] American Institute of Mathematics, Variational Methods in Celestial Mechanics; notes by Kuo-Chang Chen, version 22 June 2003, source item 21.
\url{https://aimath.org/WWN/varcelest/varcelest.pdf}.
\end{itemize}
\label{end:20002550}
\end{document}
